\section{Decomposition into \(5\) polarization modes}

Dimension five acquires tensorial content upon choosing a unit direction
of propagation \(n\) and an oriented orthonormal frame
\((e_1,e_2,n)\). Define the five real tensors
\begin{align}
E_{+}&=\frac1{\sqrt2}(e_1e_1^T-e_2e_2^T),
&E_{\times}&=\frac1{\sqrt2}(e_1e_2^T+e_2e_1^T),
\label{rm:eq:gravity-helicity2-real}\\
E_{1}&=\frac1{\sqrt2}(e_1n^T+ne_1^T),
&E_{2}&=\frac1{\sqrt2}(e_2n^T+ne_2^T),
\label{rm:eq:gravity-helicity1-real}\\
E_{0}&=\frac1{\sqrt6}(2nn^T-e_1e_1^T-e_2e_2^T).
\label{rm:eq:gravity-helicity0-real}
\end{align}
All are symmetric and traceless. The first two are transverse; the
next two mix the transverse plane with the direction of
propagation, and the last is the traceless longitudinal scalar mode.

\begin{theorem}[Complete helical decomposition of the carrier]
\label{rm:thm:gravity-five-helicities}
The tensors of
\cref{rm:eq:gravity-helicity2-real,rm:eq:gravity-helicity1-real,rm:eq:gravity-helicity0-real}
form an orthonormal basis of \(\operatorname{STF}_3\). In the
complexification, the combinations
\begin{equation}
\label{rm:eq:gravity-complex-helicities}
E_{\pm2}=\frac1{\sqrt2}(E_+\mp\ii E_\times),
\qquad
E_{\pm1}=\frac1{\sqrt2}(E_1\mp\ii E_2),
\qquad E_0=E_0
\end{equation}
are weight-\(\pm2,\pm1,0\) vectors, respectively, under rotations
about \(n\). Moreover,
\begin{equation}
\label{rm:eq:gravity-tt-on-five}
\Lambda(n)E_{\pm2}=E_{\pm2},
\qquad
\Lambda(n)E_{\pm1}=0,
\qquad
\Lambda(n)E_0=0.
\end{equation}
\end{theorem}

\begin{proof}
In the frame \((e_1,e_2,n)\), each tensor has Frobenius norm one.
Their off-diagonal entry patterns are disjoint, and the three diagonals
of \(E_+,E_0\) have zero inner product; therefore, the five tensors
are orthonormal. Since \(\dim\operatorname{STF}_3=6-1=5\), they form a basis.

A rotation through angle \(\psi\) about \(n\) sends
\((e_1,e_2)\) to
\((\cos\psi\,e_1+\sin\psi\,e_2,
-\sin\psi\,e_1+\cos\psi\,e_2)\). Substituting into the formulas gives
a rotation through angle \(2\psi\) on
\(\operatorname{span}\{E_+,E_\times\}\), one through angle \(\psi\) on
\(\operatorname{span}\{E_1,E_2\}\), and \(E_0\) fixed. Hence the weights of
\cref{rm:eq:gravity-complex-helicities}. Finally,
\(\Pi E_{+}\Pi=E_+\) and \(\Pi E_\times\Pi=E_\times\), whereas
\(\Pi E_1\Pi=\Pi E_2\Pi=0\). For \(E_0\),
\(\Pi E_0\Pi=-(e_1e_1^T+e_2e_2^T)/\sqrt6\), which is pure trace in
the plane, and subtracting \cref{rm:eq:gravity-lambda-tt} annihilates it.
\end{proof}

The theorem specifies the meaning of ``five polarizations'' without assigning all of them the same dynamic status. Before field equations are imposed, the spin-two carrier has five irreducible coordinates. In a massive Fierz--Pauli realization, the constraints leave precisely the five weights of \cref{rm:eq:gravity-complex-helicities}; in a massless realization with gauge invariance, the sectors \(\pm1\) and \(0\) are eliminated, leaving \(\pm2\). Here HMT supplies the carrier, the basis, and the projector; the choice of massive or massless dynamics belongs to the physical interface. Thus five does not mean ``five observed waves,'' but rather five tensorial degrees available before reduction.

The orientation also remains visible: replacing
\((e_1,e_2,n)\) with \((e_1,-e_2,-n)\) preserves \(E_+,E_0\) and exchanges
the phases of the helical pairs. That information is not in the
scalar \(G_a\); it resides in the carrier on which
\(\mathbb G_{5,a}\) acts.

\section{Representational chain \(12\to5\to5\to2\): carrier \(V_5\), trace-free tensor, Fierz–Pauli polarizations, and transverse helicities}

Dimension five appears in distinct objects, and only
\cref{rm:eq:gravity-12-5-5-2} authorizes transport between them:

\begin{enumerate}
\item \(V_5\) is the irreducible summand of the representation on the
twelve vertices. This is a finite result in representation theory.
\item \(\operatorname{STF}_3\) is the real five-dimensional carrier of
the \(m=-2,-1,0,1,2\) weights of spin two.
\item A massive Fierz--Pauli theory can realize these five weights
as five physical states, after fixing the action and constraints.
\item Massless general relativity retains two radiative helicities,
\(\pm2\), represented by \(E_+\) and \(E_\times\).
\item The \(E(2)\) classes of metric waves, such as \(III_5\), classify
possible amplitudes under other hypotheses; they are not automatically
\(V_5\), Fierz--Pauli, or general relativity.
\end{enumerate}

In particular, the correct statement is: \emph{HMT constructs a five-component
tensor carrier before the restrictions; the TT reduction of the massless
realization has rank two}. It is not asserted that
general relativity has five propagating polarizations.
